\documentclass[11pt]{article}
\usepackage[utf8]{inputenc}
\usepackage[T1]{fontenc}
\usepackage[french]{babel}
\usepackage[a4paper,margin=2.3cm]{geometry}
\usepackage{amsmath,amssymb}
\usepackage{booktabs}
\usepackage{graphicx}
\usepackage{caption}
\usepackage{authblk}
\usepackage{xcolor}
\usepackage[colorlinks=true,linkcolor=black,citecolor=blue,urlcolor=blue]{hyperref}

% ---- À PERSONNALISER ----
\newcommand{\Auteur}{Ronel Nkouathio}
\newcommand{\Entreprise}{[Nom de l'entreprise d'accueil]}
\newcommand{\Encadrant}{[Encadrant·e]}
% -------------------------

\title{\textbf{Mesure équitable de la fréquence cardiaque par caméra :\\
étude comparative détaillée de méthodes rPPG sur peaux foncées,\\
du fine-tuning de PhysNet aux signaux palmaires}}
\author[1,2]{\Auteur\thanks{Stage réalisé au sein de \Entreprise{}. Correspondance : nkouathioronel@gmail.com}}
\affil[1]{\Entreprise{} --- développement d'outils de santé numérique}
\affil[2]{Stage d'ingénieur / de master, encadré par \Encadrant}
\date{Juillet 2026}

\begin{document}
\maketitle

\begin{abstract}
\noindent
La photopléthysmographie à distance (rPPG) estime la fréquence cardiaque (FC)
depuis une simple vidéo. Ses performances se dégradent sur les carnations foncées
(phototypes Fitzpatrick~5--6), la mélanine absorbant le faible signal optique du
pouls --- un biais d'équité documenté, analogue à celui de l'oxymétrie de contact.
Dans le cadre d'un stage industriel visant un déploiement en pharmacie sur
smartphone, nous menons une étude comparative détaillée de méthodes rPPG sur le jeu
de données académique VitalVideos (114~sujets, majoritairement Fitzpatrick~5--6,
avec vérité-terrain PPG de contact). Nous décrivons formellement et évaluons : les
méthodes de chrominance CHROM et POS ; une variante conditionnée au teint
(CHROM-ITA) ; un réseau convolutif~1D sur signaux multi-régions ; et une stratégie
de fine-tuning de PhysNet dont nous étudions les variantes. En évaluation held-out
sans fuite, PhysNet affiné atteint une MAE de 3{,}89~bpm et le réseau~1D de
4{,}69~bpm ; couplés à une politique d'abstention par qualité de signal, l'erreur
sur les mesures acceptées tombe à 0{,}25 et 1{,}27~bpm. Nous montrons surtout que le
visage échoue sur peau foncée et proposons de mesurer sur la \textbf{paume}, peau la
moins pigmentée du corps ; un réseau~1D palmaire validé sans fuite, groupé par
personne, atteint \textbf{0{,}69~bpm de MAE sans écart entre carnations}
(Fitzpatrick~6 : 0{,}73 ; Fitzpatrick~5 : 0{,}70). Des essais préliminaires sur nos
propres vidéos smartphone confirment un signal exploitable (1--2~bpm) en capture
stable. Le code est publié en open source.

\medskip
\noindent\textbf{Mots-clés :} rPPG, fréquence cardiaque, équité, Fitzpatrick,
PhysNet, fine-tuning, CNN~1D, paume, abstention, santé numérique.
\end{abstract}

\tableofcontents
\medskip

\section{Introduction}
\label{sec:intro}

\subsection{Principe physiologique}
À chaque battement cardiaque, l'afflux de sang dans le lit capillaire cutané modifie
légèrement l'absorption de la lumière par la peau. Une caméra observe donc, sur la
peau éclairée, une infime variation périodique de couleur --- de l'ordre de $1\%$ de
l'intensité, dominante dans le canal vert (l'hémoglobine absorbant fortement autour
de 540~nm). Extraire cette composante et estimer sa fréquence donne la FC sans
contact. Le phénomène a été mis en évidence par Verkruysse
\emph{et al.}~\cite{verkruysse2008} puis automatisé par Poh \emph{et
al.}~\cite{poh2010}. Deux familles de méthodes se sont imposées : les approches
\emph{de chrominance} sans apprentissage --- CHROM~\cite{dehaan2013} et
POS~\cite{wang2017} --- et les approches d'\emph{apprentissage profond} ---
DeepPhys~\cite{chen2018}, PhysNet~\cite{yu2019}, RhythmNet~\cite{niu2020}.

\subsection{Le problème d'équité}
La quantité de \emph{mélanine} dans l'épiderme conditionne l'absorption lumineuse.
Sur une peau foncée (phototypes Fitzpatrick~5--6), une plus grande part du signal est
absorbée avant d'atteindre le lit capillaire, si bien que la modulation pulsatile
mesurable est plus faible : le rapport signal-sur-bruit (SNR) chute. Nowara
\emph{et al.}~\cite{nowara2020} le quantifient pour le rPPG. Ce biais est le pendant
optique du biais racial de l'oxymétrie de contact, établi cliniquement par Sjoding
\emph{et al.}~\cite{sjoding2020} : chez les patients à peau foncée, l'oxymètre
surestime la saturation, avec des conséquences de prise en charge. Pour un
déploiement en Afrique subsaharienne --- le cadre applicatif de ce stage --- où la
population est très majoritairement de phototype~5--6, ce biais est rédhibitoire.

\subsection{Contributions et plan}
Ce document rend compte de l'ensemble du travail de stage, sans se limiter à une
méthode. Nos contributions sont :
\begin{enumerate}\itemsep2pt
  \item une \textbf{description formelle et une comparaison rigoureuse} de cinq
  méthodes rPPG sur VitalVideos, en évaluation held-out sans fuite
  (Sections~\ref{sec:methodes}, \ref{sec:resultats}) ;
  \item une \textbf{étude de fine-tuning de PhysNet} (initialisation, gel de couches,
  entraînement de zéro) et ses enseignements (Section~\ref{sec:ft}) ;
  \item la motivation et la conception d'un \textbf{réseau convolutif~1D sur signaux
  régionaux}, léger et transférable (Section~\ref{sec:cnn1d}) ;
  \item une \textbf{politique d'abstention} à quatre garde-fous, pour ne rendre
  qu'une valeur fiable (Section~\ref{sec:abstention}) ;
  \item notre contribution principale : mesurer sur la \textbf{paume}, avec un
  réseau~1D palmaire validé sans fuite qui \emph{supprime l'écart d'équité}
  (Section~\ref{sec:paume}) ;
  \item des \textbf{tests préliminaires smartphone} et un \textbf{rapport honnête des
  échecs} (Sections~\ref{sec:perso}, \ref{sec:negatifs}).
\end{enumerate}

\paragraph{Positionnement et humilité.} Ce travail est un travail de \emph{stage},
exploratoire et industriel, non une étude clinique. Les résultats sont des preuves de
concept. Nous privilégions la rigueur méthodologique (absence de fuite) et la
transparence sur les limites et les échecs.

\section{Cadres de travail}
\label{sec:cadres}

\paragraph{Cadre académique.} Les algorithmes de référence et les jeux de données
publics forment le socle scientifique. Nous validons d'abord sur
UBFC-rPPG~\cite{bobbia2019}, benchmark standard, puis principalement sur
\textbf{VitalVideos}~\cite{toye2023}.

\paragraph{Cadre industriel.} L'objectif est un outil utilisable en pharmacie, sur
smartphone, par des non-spécialistes. Cela impose des contraintes absentes des
benchmarks : compression H.264, réglages caméra variables, éclairage non contrôlé,
et une population cible à peau foncée. Ces contraintes orientent nos choix (site
palmaire, abstention, protocole de capture).

\paragraph{Cadre open source.} Nous nous appuyons sur rPPG-Toolbox~\cite{liu2023}
(implémentations de référence et poids pré-entraînés), MediaPipe~\cite{lugaresi2019}
(détection visage/main) et un \emph{face parsing} BiSeNet~\cite{yu2018bisenet}. En
retour, le code de nos pipelines et évaluations est publié en open source.

\section{Données et prétraitement}
\label{sec:donnees}

\subsection{Le jeu de données VitalVideos}
Chaque sujet de VitalVideos est filmé selon plusieurs \emph{scénarios} :
\emph{faceonly} (visage seul), \emph{facepalm} (visage et paume dans le cadre),
\emph{handheld} (caméra tenue en main), \emph{exertion} (après effort),
\emph{slowbreath}/\emph{quickbreath} (respiration contrôlée). Pour chaque scénario, le
JSON associé fournit : la vidéo RGB (résolution $\sim$$960\times1056$, 60~fps,
encodage quasi sans perte), une PPG de contact synchronisée (oxymètre \emph{CMS50},
colonnes temps/PPG/FC/SpO$_2$), une pression artérielle au brassard (Omron), un
signal de pression respiratoire, et des métadonnées (âge, genre, phototype de
Fitzpatrick, lieu, environnement). Cette richesse, et surtout la forte proportion de
phototypes~5--6, en font un support rare pour l'étude de l'équité.

\subsection{Alignement de la vérité-terrain}
La FC de référence est dérivée de la PPG de contact : la colonne PPG est filtrée
passe-bande à la fréquence native du capteur, puis interpolée sur les
\emph{horodatages vidéo} (fournis image par image dans le JSON). Toutes les vidéos
sont ramenées à 30~images/s. Ce rééchantillonnage est fait \emph{uniquement vers le
bas} (décimation depuis 60~fps) : sur-échantillonner (p.\,ex. $24\to30$) dupliquerait
des images et détruirait la représentation \emph{DiffNormalized} de PhysNet.

\subsection{Deux représentations d'entrée}
Selon la méthode, une vidéo est transformée en :

\paragraph{(a) Clips normalisés (PhysNet).} Le visage est suivi et recadré, chaque
image redimensionnée en $72\times72$, puis convertie en \emph{DiffNormalized} : la
différence normalisée trame à trame
\begin{equation}
d(t) = \frac{I(t+1) - I(t)}{I(t+1) + I(t) + \epsilon},
\end{equation}
qui accentue la composante pulsatile et atténue l'éclairage stationnaire. Les clips
font 128~images.

\paragraph{(b) Signaux multi-régions (CHROM/POS/CNN~1D).} Pour chaque image, on
calcule la couleur moyenne de régions cutanées. Pour le \emph{visage},
$N_r = 23$~régions : six régions anatomiques (front, deux joues, nez, deux
sous-yeux) délimitées par des points-clés MediaPipe FaceMesh, la peau complète, et
une grille $4\times4$ (16~cellules) couvrant la boîte du visage, à la manière des
cartes spatio-temporelles MSTmap~\cite{niu2020}. Un masque de peau produit par le
\emph{face parsing} BiSeNet exclut cheveux, yeux et arrière-plan. Chaque région est
moyennée dans \emph{trois espaces colorimétriques} --- RGB, YUV, Lab --- soit
$N_c = 9$~canaux. On obtient un tenseur $x \in \mathbb{R}^{T \times N_r \times N_c}$
($23\times9 = 207$~canaux pour le visage). Les images sans détection sont interpolées
linéairement.

\subsection{Filtrage et held-out}
Tous les signaux estimés sont filtrés passe-bande dans $[0{,}7;2{,}5]$~Hz
($\approx$$[42;150]$~bpm). Vingt-deux sujets ($\sim$54~scénarios) sont réservés au
test et n'apparaissent jamais à l'entraînement (Section~\ref{sec:rigueur}).

\section{Méthodes évaluées}
\label{sec:methodes}

\subsection{Normalisation temporelle}
Toutes les méthodes de chrominance et le CNN~1D partent d'une normalisation
temporelle par canal, qui retire la composante continue (teint, éclairage) et ne
garde que les fluctuations relatives :
\begin{equation}
\tilde{x}_c(t) = \frac{x_c(t)}{\bar{x}_c} - 1, \qquad
\bar{x}_c = \frac{1}{T}\sum_t x_c(t).
\end{equation}

\subsection{CHROM}
CHROM~\cite{dehaan2013} construit deux signaux de chrominance à partir des canaux RGB
normalisés :
\begin{equation}
X = 3\tilde{R} - 2\tilde{G}, \qquad Y = 1{,}5\tilde{R} + \tilde{G} - 1{,}5\tilde{B},
\end{equation}
puis combine $S = X - \alpha Y$ avec $\alpha = \sigma(X)/\sigma(Y)$, ce qui annule au
premier ordre la composante de mouvement (spéculaire) sous l'hypothèse d'un teint
standardisé. Aucun entraînement n'est requis.

\subsection{POS}
POS~\cite{wang2017} projette les RGB temporellement normalisés sur un plan orthogonal
à la direction de l'intensité :
\begin{equation}
\begin{bmatrix} S_1 \\ S_2 \end{bmatrix} =
\begin{bmatrix} 0 & 1 & -1 \\ -2 & 1 & 1 \end{bmatrix}
\begin{bmatrix} \tilde{R} \\ \tilde{G} \\ \tilde{B} \end{bmatrix},
\qquad S = S_1 + \frac{\sigma(S_1)}{\sigma(S_2)} S_2.
\end{equation}
POS est en général plus robuste que CHROM sur peaux claires ; nous vérifions son
comportement sur peaux foncées.

\subsection{CHROM conditionné au teint : CHROM-ITA}
\label{sec:chromita}
La projection de CHROM est optimale pour un teint \og{}moyen\fg{}. Nous conditionnons ses
coefficients au teint via l'\emph{Individual Typology Angle} (ITA), descripteur
colorimétrique standard en dermatologie :
\begin{equation}
\mathrm{ITA} = \arctan\!\left(\frac{L^* - 50}{b^*}\right)\cdot\frac{180}{\pi},
\end{equation}
où $(L^*, b^*)$ sont mesurés dans l'espace CIE-Lab sur la peau (ITA élevé $=$ peau
claire, ITA fortement négatif $=$ peau foncée). Un petit perceptron
$\phi_\theta : \mathrm{ITA} \mapsto (a_1,\dots,a_k)$ prédit les coefficients de la
combinaison chrominance, permettant d'\emph{adapter} la projection à chaque
carnation. Ce module est entraîné sur VitalVideos. Pour une estimation d'ITA robuste
à l'éclairage, nous normalisons par la sclère (blanc de l'\oe{}il) comme référence
d'illumination.

\subsection{Signaux multi-régions et réseau convolutif~1D}
\label{sec:cnn1d}

\paragraph{Motivation.} PhysNet opère par convolutions 3D sur la vidéo brute :
puissant mais coûteux (GPU, plein cadre) et sensible aux artefacts de compression.
Nous proposons une voie plus légère : \emph{réduire la vidéo à un ensemble compact de
signaux régionaux multi-colorimétriques} (Section~\ref{sec:donnees}), puis régresser
le pouls par un réseau \textbf{convolutif~1D dilaté}. Trois motivations : (i)~\emph{
coût} très réduit ($\sim$110\,000~paramètres, pas de convolution 3D), exécutable sur
matériel modeste ; (ii)~\emph{robustesse} par la diversité spatiale (multi-régions,
certaines mieux perfusées) et colorimétrique (RGB+YUV+Lab, l'illumination affectant
différemment chaque espace) ; (iii)~\emph{transférabilité} : la même représentation
s'applique directement à la \emph{paume}, pour laquelle il n'existe pas de modèle de
segmentation dédié --- ce choix est le pont vers notre contribution principale
(Section~\ref{sec:paume}). Conceptuellement, le CNN~1D est la version \emph{apprenable}
du pipeline CHROM/POS multi-ROI : il apprend implicitement la pondération des régions,
la projection chrominance et le filtrage temporel.

\paragraph{Architecture.} L'entrée $x \in \mathbb{R}^{B\times C\times T}$ ($C = N_r
N_c$) passe par : une \emph{tête d'entrée} $1\times1$ ($C \to 64$ canaux) avec
BatchNorm et ReLU ; puis quatre \emph{blocs temporels résiduels} de dilatations
$\{1,2,4,8\}$, chacun composé de deux convolutions 1D (noyau~3, canaux~64) avec
BatchNorm, ReLU et Dropout~$0{,}1$, et connexion résiduelle ; enfin une \emph{tête de
projection} $1\times1$ ($64 \to 1$). Les dilatations croissantes donnent un champ
réceptif temporel d'environ 60~échantillons ($\sim$2~s à 30~fps) pour un coût
paramétrique faible. La sortie est l'onde PPG estimée $\hat{y} \in \mathbb{R}^{B\times
T}$.

\paragraph{Perte.} L'apprentissage minimise la \emph{corrélation de Pearson négative}
entre l'onde prédite et l'onde de référence, invariante à l'amplitude et au décalage :
\begin{equation}
\mathcal{L}(\hat{y}, y) = -\,
\frac{\sum_t (\hat{y}_t - \bar{\hat y})(y_t - \bar y)}
{\sqrt{\sum_t (\hat{y}_t - \bar{\hat y})^2}\,\sqrt{\sum_t (y_t - \bar y)^2} + \epsilon}.
\end{equation}
La FC est ensuite l'argument du pic dominant du périodogramme dans la bande
cardiaque.

\subsection{PhysNet et fine-tuning}
\label{sec:physnet}
PhysNet~\cite{yu2019} est un réseau encodeur-décodeur 3D (convolutions
spatio-temporelles, \emph{max-pooling} temporel puis \emph{upsampling}) qui régresse
l'onde PPG depuis les clips DiffNormalized. Des poids pré-entraînés existent (PURE,
SCAMPS, UBFC). Le sujet central est l'\emph{adaptation} de ces poids à VitalVideos ;
nous en faisons une étude dédiée (Section~\ref{sec:ft}).

\subsection{Ballistocardiographie (rBCG)}
En complément \emph{optique}, nous estimons la FC par le micro-mouvement de la tête
induit par l'éjection sanguine~\cite{balakrishnan2013} : suivi de points rigides par
flux optique de Lucas-Kanade, sélection des trajectoires stables, filtrage dans la
bande cardiaque, décomposition en composantes principales (PCA/SVD), et choix de la
composante la plus périodique (SNR maximal). Étant \emph{mécanique}, le rBCG est
\emph{indépendant de la mélanine} --- atout sur peau foncée --- mais bruité et
sensible au mouvement ; nous l'utilisons comme \emph{votant}, jamais comme mesure
autonome.

\subsection{Fusion adaptative et abstention}
\label{sec:abstention}
\paragraph{Fusion.} Les FC des différentes méthodes sont agrégées par une règle
adaptative : \emph{médiane} si les méthodes s'accordent (faible écart-type),
\emph{sélection au SNR} sinon, avec un garde-fou anti-harmonique pour éviter de
choisir un $\times$2 fréquentiel.

\paragraph{Quatre garde-fous.} Inspirés des systèmes commerciaux, nous ne rendons une
valeur que si la mesure est jugée \emph{fiable}. Une mesure est \textsc{fiable} si :
(i)~\emph{accord} --- l'écart-type des FC est faible ; (ii)~\emph{SNR} médian
suffisant ; (iii)~\emph{couverture} --- une analyse de qualité par fenêtres
glissantes (fenêtre 10~s, pas 2~s) juge chaque fenêtre et l'on agrège les fenêtres
fiables, la fraction fiable donnant la couverture ; (iv)~\emph{absence d'ambiguïté}
--- deux rythmes candidats de puissance spectrale comparable (le second $\geq 60\%$ du
premier, séparés de plus de 4~bpm, non harmoniques) déclenchent un doute. Faute de
quoi, la sortie est \og{}à confirmer\fg{} ou \og{}à refaire\fg{}. Les seuils exacts sont donnés en
annexe~\ref{app:seuils}.

\paragraph{Qualité de signal (SNR).} Le SNR aveugle mesure la concentration
spectrale autour de la FC estimée et de sa première harmonique :
\begin{equation}
\mathrm{SNR} = 10\log_{10}
\frac{\sum_{f \in \mathcal{B}(f_{HR})} P(f)}{\sum_{f \notin \mathcal{B}(f_{HR})} P(f)},
\end{equation}
où $\mathcal{B}$ est une bande étroite autour de $f_{HR}$ et $2f_{HR}$. C'est
l'indicateur pivot de la politique d'abstention.

\section{Étude de fine-tuning de PhysNet}
\label{sec:ft}

Partant d'un PhysNet pré-entraîné, quelle stratégie d'adaptation à VitalVideos
généralise le mieux ? Nous comparons (Table~\ref{tab:ft}) : (A)~initialisation
pré-entraînée $+$ fine-tuning \emph{complet} (toutes couches) ; (B, C)~\emph{gel} de
l'encodeur (seul le décodeur réentraîné, deux variantes) ; (D)~entraînement de
\emph{zéro} (init.\ aléatoire) ; (E)~un entraînement sur $\sim$100~sujets, sur un
split de test différent. L'optimiseur est Adam (pas $10^{-3}$, \emph{weight decay}
$10^{-4}$), 30~époques, la meilleure époque étant retenue sur la validation.

\begin{table}[h]
\centering
\caption{Étude de fine-tuning de PhysNet (validation du fine-tuning). Le modèle
retenu (A) est celui évalué en comparaison held-out (Table~\ref{tab:visage}).}
\label{tab:ft}
\begin{tabular}{llc}
\toprule
 & Stratégie & MAE (bpm) \\
\midrule
A & Init.\ pré-entraînée $+$ fine-tuning \textbf{complet} & \textbf{2{,}1} \\
D & Entraînement de zéro (aléatoire) & 5{,}5 \\
B & Gel de l'encodeur (variante 1) & 21{,}7 \\
C & Gel de l'encodeur (variante 2) & 23{,}0 \\
E & $\sim$100 sujets, split différent (confondu) & 14{,}7 \\
\bottomrule
\end{tabular}
\end{table}

\paragraph{Enseignements.} (1)~Le \emph{gel} de l'encodeur est catastrophique
(21--23~bpm) : les features 3D pré-entraînées sur d'autres domaines (visages clairs,
données synthétiques) ne se transposent pas telles quelles ; il faut réentraîner tout
le réseau. (2)~L'\emph{initialisation pré-entraînée aide nettement} (2{,}1 contre
5{,}5 de zéro, soit un facteur $\sim$2{,}7) : le pré-entraînement fournit un a priori
utile malgré le changement de domaine. (3)~L'expérience à 100~sujets, sur un split de
test différent et des extractions hétérogènes, est \emph{non concluante} (confondue)
--- nous ne surinterprétons pas ce chiffre. Nous retenons la stratégie~A.

\section{Protocole d'évaluation et rigueur}
\label{sec:rigueur}

\paragraph{Absence de fuite.} Un modèle profond peut \emph{mémoriser} un sujet (teint,
morphologie, éclairage) ; l'évaluer sur des fenêtres du même sujet qu'à
l'entraînement gonfle artificiellement les scores. Nous imposons une séparation
\textbf{au niveau de la personne} : toutes les prises d'un individu sont entièrement
en apprentissage \emph{ou} en test. Pour la validation palmaire, nous avons constaté
que quatre personnes figuraient sous deux identifiants de dossier distincts ; nous
fusionnons ces dossiers (union par identifiant de session) avant la partition.

\paragraph{Une leçon de fuite annexe.} Nous avons vérifié empiriquement qu'entraîner
le module d'abstention sur des sujets déjà vus par les modèles \emph{dégrade} le
résultat held-out (la MAE passe de $1{,}15$ à $2{,}48$~bpm) : les SNR de ces sujets
mémorisés sont artificiellement élevés, si bien que le module apprend à leur faire
trop confiance. Le module de qualité doit donc être calibré sur des sujets
\emph{jamais vus} par les modèles.

\paragraph{Métriques.} Nous rapportons la MAE de FC ($\frac1N\sum|\hat{f}-f|$), le
pourcentage de mesures à moins de 5~bpm, le SNR aveugle, et la corrélation de Pearson
d'onde. L'évaluation visage est au niveau scénario sur 22~sujets ; la validation
palmaire est une \textbf{validation croisée à 5~plis, groupée par personne} sur
108~individus, un modèle réentraîné de zéro par pli et les prédictions held-out
agrégées.

\section{Résultats}
\label{sec:resultats}

\subsection{Comparaison des méthodes sur VitalVideos (visage)}
La Table~\ref{tab:visage} rassemble les cinq méthodes sur le held-out (54~scénarios,
22~sujets jamais vus). PhysNet affiné domine en brut (MAE~3{,}89~bpm, 83\% à moins de
5~bpm), suivi du réseau~1D (4{,}69). Les méthodes sans modèle sont nettement derrière
(CHROM~8{,}69 ; POS~9{,}90), CHROM-ITA améliorant légèrement CHROM (8{,}27), ce qui
confirme l'intérêt --- modeste --- du conditionnement au teint.

\begin{table}[h]
\centering
\caption{Comparaison held-out sur VitalVideos (54~scénarios, 22~sujets sans fuite).
La colonne \og{}validés\fg{} applique la politique d'abstention (on ne retient que les
mesures de SNR $>-1$~dB) : l'erreur sur les mesures \emph{acceptées} chute
fortement.}
\label{tab:visage}
\begin{tabular}{lccccc}
\toprule
Méthode & MAE & \%$<$5 & Validés & MAE (validés) \\
\midrule
PhysNet (fine-tuné, A) & \textbf{3{,}89} & 83\% & 35/52 & \textbf{0{,}25} \\
CNN~1D (régions)       & 4{,}69 & 80\% & 40/54 & 1{,}27 \\
CHROM-ITA              & 8{,}27 & 67\% & 21/54 & 4{,}27 \\
CHROM                  & 8{,}69 & 59\% & 17/54 & 3{,}52 \\
POS                    & 9{,}90 & 65\% & 23/54 & 3{,}90 \\
\midrule
Fusion adaptative      & 4{,}69 & 81\% & \multicolumn{2}{c}{consensus $\times43$, sélection $\times11$} \\
\bottomrule
\end{tabular}
\end{table}

\paragraph{L'apport de l'abstention.} La colonne \og{}MAE (validés)\fg{} est éloquente : en
ne conservant que les mesures de bon SNR, PhysNet passe de 3{,}89 à
\textbf{0{,}25~bpm} et le réseau~1D de 4{,}69 à 1{,}27~bpm. Le prix est une
\emph{couverture} partielle (35/52 et 40/54). Le réseau~1D offre la meilleure
couverture (40/54) ; PhysNet la meilleure précision sur mesures acceptées. Ce
compromis \emph{couverture / précision} est au c\oe{}ur de l'approche : en contexte
médical, mieux vaut s'abstenir que rendre une valeur fausse. La fusion adaptative
choisit le consensus dans 43~cas sur 54 et bascule en sélection-SNR dans les
11~cas de désaccord.

\subsection{Le constat d'inéquité et la piste palmaire}
\label{sec:paume}
Le point critique apparaît en stratifiant par phototype : sur peau très foncée, le
\emph{visage} devient peu fiable. Nous proposons de mesurer sur la \textbf{paume},
seule peau glabre faiblement pigmentée quel que soit le phototype. Le pipeline
palmaire reprend la Section~\ref{sec:cnn1d} : 18~régions issues de MediaPipe~Hands
(huit régions anatomiques --- centre, thénar, hypothénar, sous-doigts, poignet ---,
une grille $3\times3$, la paume entière), $\times 9$~canaux, soit 162~canaux
d'entrée. Le jeu palmaire compte 127~enregistrements de 108~personnes (75~Fitz\,6,
48~Fitz\,5, 4~Fitz\,4), extraits des scénarios \emph{facepalm} et \emph{handheld}.

\paragraph{Paume contre visage (méthodes classiques).} Sur 24~vidéos, la paume
surclasse nettement le visage sur peau foncée (Table~\ref{tab:palmvsface}) : sur
Fitzpatrick~6, le visage s'effondre (MAE~12{,}4, SNR négatif) tandis que la paume
reste précise (MAE~1{,}1, SNR~$+7{,}1$~dB). Le SNR palmaire est \emph{positif partout},
le SNR facial \emph{négatif partout}.

\begin{table}[h]
\centering
\caption{Paume contre visage (meilleure méthode sans modèle par site), par phototype.
MAE en bpm, SNR aveugle en dB.}
\label{tab:palmvsface}
\begin{tabular}{lcccc}
\toprule
 & \multicolumn{2}{c}{\textbf{Paume}} & \multicolumn{2}{c}{\textbf{Visage}} \\
\cmidrule(lr){2-3}\cmidrule(lr){4-5}
Phototype & MAE & SNR & MAE & SNR \\
\midrule
Fitzpatrick 5 ($n=7$)  & 1{,}3 & $+5{,}9$ & 2{,}4  & $-1{,}6$ \\
Fitzpatrick 6 ($n=17$) & 1{,}1 & $+7{,}1$ & 12{,}4 & $-2{,}5$ \\
\midrule
Global ($n=24$)        & 1{,}2 & $+6{,}7$ & 9{,}5  & $-2{,}2$ \\
\bottomrule
\end{tabular}
\end{table}

\paragraph{Réseau palmaire (validation croisée sans fuite).} Le réseau~1D palmaire,
évalué par validation croisée groupée par personne (Table~\ref{tab:cv}), atteint une
MAE globale de \textbf{0{,}69~bpm} (98\% à moins de 5~bpm, $r=0{,}88$). Surtout,
\textbf{l'écart entre carnations disparaît} : Fitzpatrick~6 (0{,}73) est aussi bon que
Fitzpatrick~5 (0{,}70). Les cinq plis donnent des MAE de $1{,}76 / 0{,}14 / 0{,}06 /
0{,}24 / 1{,}38$~bpm.

\begin{table}[h]
\centering
\caption{Réseau~1D palmaire --- validation croisée 5~plis, groupée par personne, sans
fuite (127~enregistrements, 108~personnes).}
\label{tab:cv}
\begin{tabular}{lcccc}
\toprule
Phototype & $n$ & MAE (bpm) & SNR (dB) & $r$ (onde) \\
\midrule
Fitzpatrick 4 & 4  & 0{,}00 & $+8{,}1$ & 0{,}90 \\
Fitzpatrick 5 & 48 & 0{,}70 & $+6{,}8$ & 0{,}89 \\
Fitzpatrick 6 & 75 & 0{,}73 & $+6{,}7$ & 0{,}88 \\
\midrule
Global        & 127 & \textbf{0{,}69} & $+6{,}8$ & 0{,}88 \\
\bottomrule
\end{tabular}
\end{table}

\paragraph{Prudence d'interprétation.} Sur un signal palmaire propre de $\sim$20~s,
estimer la \emph{fréquence} dominante est une tâche \og{}facile\fg{} ; la MAE très basse en
rend partiellement compte. La corrélation d'onde ($r=0{,}88$ sur sujets jamais vus)
est l'indicateur le plus honnête que le modèle apprend la \emph{forme} du pouls, non
seulement sa fréquence. Une expérience de contrôle montre par ailleurs que réduire
l'entrée aux seuls canaux RGB (54~canaux) donne des résultats identiques : les canaux
YUV/Lab sont redondants pour ce site.

\subsection{Tests sur nos propres vidéos (smartphone)}
\label{sec:perso}
Nous avons filmé nos propres vidéos avec l'application open source \emph{Open Camera},
sans instrumentation de référence (analyse qualitative). La Table~\ref{tab:phone}
résume cinq prises. Trois enseignements se dégagent.

\begin{table}[h]
\centering
\caption{Prises smartphone (paume). La qualité (SNR, verdict) ne suit pas le débit
d'encodage mais la stabilité et l'éclairage de la prise.}
\label{tab:phone}
\begin{tabular}{lccccc}
\toprule
Prise & cadrage & fps & débit & SNR paume & verdict \\
\midrule
1 & paume seule    & 24{,}5 & 39~Mbit/s & $+2{,}9$ & GO (err 1) \\
2 & main+visage    & 29{,}2 & 56~Mbit/s & $+2{,}0$ & GO (err 2) \\
3 & main+visage    & 29{,}6 & 40~Mbit/s & $-2{,}3$ & refaire \\
4 & main+visage    & 29{,}3 & 60~Mbit/s & $-4{,}0$ & refaire \\
5 & main+visage    & 19{,}7 & 59~Mbit/s & $-1{,}4$ & refaire \\
\bottomrule
\end{tabular}
\end{table}

(1)~Sur une \emph{bonne} prise (paume posée, stable, éclairée), le réseau palmaire
retrouve la FC à \textbf{1--2~bpm} près, avec SNR positif et accord inter-méthodes
(Fig.~\ref{fig:onde}). (2)~Le facteur dominant de qualité n'est \emph{pas} le débit
d'encodage --- une prise à 60~Mbit/s échoue et une à 39~Mbit/s réussit --- mais la
\textbf{stabilité et l'éclairage} ; sur les prises dégradées, le système
\emph{s'abstient}, comme attendu. (3)~Sur ces prises smartphone à peau foncée, le
\emph{visage} échoue (PhysNet jusqu'à 20~bpm d'erreur) là où la paume réussit,
confirmant qualitativement le résultat d'équité en conditions réelles.

\begin{figure}[h]
\centering
\includegraphics[width=0.86\linewidth]{figures/onde_paume_smartphone.png}
\caption{Onde rPPG palmaire estimée sur une vidéo smartphone. Le système déclare la
prise fiable (\og{}GO\fg{}) ; les trois méthodes s'accordent et le spectre présente un pic
net, malgré la compression H.264 et une capture non instrumentée.}
\label{fig:onde}
\end{figure}

\subsection{Résultats négatifs (rapportés honnêtement)}
\label{sec:negatifs}
Pour éviter à d'autres des impasses coûteuses :
\begin{itemize}\itemsep2pt
  \item \textbf{Apprentissage auto-supervisé} (styles SiNC~\cite{speth2023},
  Contrast-Phys~\cite{sun2022}) : nos implémentations sur signaux régionaux se sont
  effondrées (sortie constante) ou n'ont pas appris le pouls (corrélation nulle).
  Ces méthodes exigent la vidéo brute multi-signaux et un réglage fin que notre
  représentation ne permet pas.
  \item \textbf{Pression artérielle} par morphologie du pouls : indissociable d'un
  simple biais d'âge (AUC au niveau du hasard, $\approx0{,}50$). La voie du temps de
  transit (deux sites optiques, p.\,ex.\ paume et visage, annulant la période de
  pré-éjection) est théoriquement plus fondée mais bute sur la précision temporelle à
  30--60~fps et nécessite une calibration par sujet.
  \item \textbf{Variabilité fine de la FC} (RMSSD) : hors de portée, limitée par la
  résolution temporelle de la caméra (33~ms à 30~fps) et l'arrondi de l'onde rPPG. La
  variabilité \emph{grossière} (SDNN, LF/HF) reste envisageable avec haut débit
  d'images et fusion paume$+$visage.
\end{itemize}

\section{Discussion}
\label{sec:discussion}

Le résultat central est double. D'une part, sur le visage, l'apprentissage profond
(PhysNet, CNN~1D) domine nettement les méthodes de chrominance sur VitalVideos, et le
fine-tuning \emph{complet} depuis une initialisation pré-entraînée est la stratégie
gagnante. D'autre part, aucune méthode faciale ne résout le biais d'équité : sur peau
très foncée, le signal facial est intrinsèquement faible. Le déplacement vers la
\emph{paume} contourne ce mur physiologique, car la peau palmaire est faiblement
pigmentée indépendamment du phototype. C'est, à notre connaissance, une piste peu
exploitée dans la littérature rPPG, largement centrée sur le visage pour des raisons
d'usage (surveillance passive) et de disponibilité de données. Notre cadre
applicatif --- point de soin coopératif, population à peau foncée, jeu de données
riche en \emph{facepalm} --- réunit précisément les conditions où cette piste devient
pertinente.

\section{Limites}
\label{sec:limites}
(1)~Le résultat palmaire est obtenu sur caméra industrielle (60~fps, quasi sans
perte) ; la validation smartphone reste préliminaire ($n$ faible, sans vérité-terrain
instrumentée). (2)~Seule la FC est étudiée ; SpO$_2$, pression et respiration ne sont
pas atteignables de façon fiable avec ces données. (3)~L'échantillon Fitzpatrick~4 est
très petit. (4)~Aucune validation clinique ni test utilisateur. Le statut est celui
d'une \emph{preuve de concept forte}, non d'un produit validé.

\section{Travaux futurs}
\label{sec:futur}
La priorité est une \textbf{collecte smartphone contrôlée} : 15--20~sujets, dont
$\geq$10 de phototype~5--6, paume posée et stabilisée, éclairage constant, et une FC
de référence par oxymètre de contact. Elle permettra (i)~d'établir statistiquement le
transfert de l'avantage palmaire au smartphone, (ii)~d'affiner le modèle sur le
domaine cible, (iii)~de calibrer l'abstention en conditions réelles. Les outils
(contrôle qualité en direct d'une prise, journalisation de la référence, analyse
batch stratifiée) sont préparés et publiés.

\section{Conclusion}
\label{sec:conclusion}
Ce travail de stage présente une étude comparative détaillée de méthodes rPPG sur
peaux foncées. Sur le visage, PhysNet affiné (MAE~3{,}89, 0{,}25 sur mesures
acceptées) devance le réseau~1D et les méthodes classiques ; l'étude de fine-tuning
montre qu'il faut réentraîner tout le réseau depuis une initialisation
pré-entraînée. Mais le visage reste peu fiable sur peau très foncée ; notre
contribution est d'y répondre en mesurant sur la \textbf{paume}, où un réseau~1D
validé sans fuite atteint 0{,}69~bpm de MAE \emph{sans écart entre carnations}. Nous
insistons sur la rigueur (validation sans fuite), l'honnêteté (échecs rapportés) et
le caractère open source de la démarche.

\section*{Remerciements}
L'auteur remercie \Entreprise{} pour l'accueil et l'encadrement de ce stage, ainsi
que les auteurs de VitalVideos et des bibliothèques open source (rPPG-Toolbox,
MediaPipe, BiSeNet). Ce document rend compte d'un travail d'étudiant, avec les
limites que cela suppose.

\appendix
\section{Seuils de la politique d'abstention}
\label{app:seuils}
Accord : \textsc{fiable} si écart-type $\leq 6$~bpm, \textsc{rejet} si $>12$~bpm.
SNR médian : \textsc{fiable} si $>-3{,}5$~dB, \textsc{rejet} si $\leq -4{,}5$~dB.
Témoin isolé : une méthode de SNR $>+1$~dB, détachée d'au moins 2~dB de la suivante,
non harmonique, peut valider seule (le rBCG en est exclu, son SNR n'étant pas à la
même échelle que l'optique). Couverture SQA : fenêtre 10~s, pas 2~s ; \textsc{fiable}
requiert une couverture $\geq 0{,}20$, \textsc{rejet} si $<0{,}05$. Ambiguïté : le
deuxième pic spectral (maximum local, séparé de $>4$~bpm, non harmonique) déclenche un
doute s'il atteint $\geq 60\%$ du premier.

\section{Hyperparamètres}
\label{app:hp}
CNN~1D : 64~canaux cachés, dilatations $\{1,2,4,8\}$, Dropout~$0{,}1$,
$\sim$110\,000~paramètres. Fenêtres de $\mathrm{CLIP\_LEN}=128$~images
($\approx$4{,}3~s à 30~fps). Optimiseur Adam, pas $10^{-3}$, \emph{weight decay}
$10^{-4}$, ordonnancement cosinus, 60~époques, taille de lot~16. PhysNet : clips de
128~images, entrée DiffNormalized $72\times72$, fine-tuning 30~époques. Bande
d'analyse : $[0{,}7;2{,}5]$~Hz.

\begin{thebibliography}{99}
\small
\bibitem{verkruysse2008} W.~Verkruysse, L.~O.~Svaasand, J.~S.~Nelson.
  \emph{Remote plethysmographic imaging using ambient light.} Optics Express,
  16(26):21434--21445, 2008.
\bibitem{poh2010} M.-Z.~Poh, D.~J.~McDuff, R.~W.~Picard.
  \emph{Non-contact, automated cardiac pulse measurements using video imaging and
  blind source separation.} Optics Express, 18(10):10762--10774, 2010.
\bibitem{dehaan2013} G.~de~Haan, V.~Jeanne.
  \emph{Robust pulse rate from chrominance-based rPPG.} IEEE Trans. Biomedical
  Engineering, 60(10):2878--2886, 2013.
\bibitem{wang2017} W.~Wang, A.~C.~den~Brinker, S.~Stuijk, G.~de~Haan.
  \emph{Algorithmic principles of remote PPG.} IEEE Trans. Biomedical Engineering,
  64(7):1479--1491, 2017.
\bibitem{chen2018} W.~Chen, D.~McDuff.
  \emph{DeepPhys: Video-based physiological measurement using convolutional
  attention networks.} ECCV, 2018.
\bibitem{yu2019} Z.~Yu, X.~Li, G.~Zhao.
  \emph{Remote photoplethysmograph signal measurement from facial videos using
  spatio-temporal networks (PhysNet).} BMVC, 2019.
\bibitem{niu2020} X.~Niu \emph{et al.}
  \emph{RhythmNet: End-to-end heart rate estimation from face via spatial-temporal
  representation.} IEEE Trans. Image Processing, 29:2409--2423, 2020.
\bibitem{balakrishnan2013} G.~Balakrishnan, F.~Durand, J.~Guttag.
  \emph{Detecting pulse from head motions in video.} CVPR, 2013.
\bibitem{nowara2020} E.~M.~Nowara, D.~McDuff, A.~Veeraraghavan.
  \emph{A meta-analysis of the impact of skin type and gender on non-contact
  photoplethysmography measurements.} CVPR Workshops, 2020.
\bibitem{sjoding2020} M.~W.~Sjoding \emph{et al.}
  \emph{Racial bias in pulse oximetry measurement.} New England Journal of Medicine,
  383:2477--2478, 2020.
\bibitem{speth2023} J.~Speth \emph{et al.}
  \emph{Non-contrastive unsupervised learning of physiological signals from video
  (SiNC).} CVPR, 2023.
\bibitem{sun2022} Z.~Sun, X.~Li.
  \emph{Contrast-Phys: Unsupervised video-based remote physiological measurement via
  spatiotemporal contrast.} ECCV, 2022.
\bibitem{toye2023} P.-J.~Toye.
  \emph{VitalVideos: A large-scale dataset of face videos with PPG and blood pressure
  ground truth.} arXiv:2306.11891, 2023.
\bibitem{bobbia2019} S.~Bobbia \emph{et al.}
  \emph{Unsupervised skin tissue segmentation for remote photoplethysmography
  (UBFC-rPPG).} Pattern Recognition Letters, 124:82--90, 2019.
\bibitem{liu2023} X.~Liu \emph{et al.}
  \emph{rPPG-Toolbox: Deep remote PPG toolbox.} NeurIPS Datasets and Benchmarks, 2023.
\bibitem{lugaresi2019} C.~Lugaresi \emph{et al.}
  \emph{MediaPipe: A framework for building perception pipelines.} arXiv:1906.08172,
  2019.
\bibitem{yu2018bisenet} C.~Yu \emph{et al.}
  \emph{BiSeNet: Bilateral segmentation network for real-time semantic segmentation.}
  ECCV, 2018.
\end{thebibliography}

\end{document}
